\subsection{Learning to Produce Artifacts Without External Skills}
\label{sec:internalization}

We further ask whether Skill-informed demonstrations can improve a small
model when the procedural documents are absent at inference. Inspired by
SkillGym's distinction between external assistance and parameter transfer
\citep{skillgym}, we compare Qwen3-4B-Instruct-2507, LoRA-SFT, and subsequent
outcome-based GRPO on the same 305 tasks under No Skill.

\paragraph{Teach decisions through complete artifacts.}
We decompose the 14 Skills into operations and evidence that should change
their application: repeated failed hints, for example, should trigger a
different response from an initial request for help. EduBench question seeds
are reconstructed around complete-product creation, changed learner evidence,
resource boundaries, and artifact repair. Screening and revision retain
246 demonstrations: 162 generated reference artifacts, 52 revised artifacts,
and 32 final answers from environment executions. Each student input contains
the task; the supervised target is its complete teaching artifact. Skill
documents and rubrics are omitted from student inputs. This supervision
targets educational decisions expressed in products, rather than native
tool-call trajectories.

SFT applies completion-only loss for one epoch, preserving all complete
targets within 8,192 tokens. GRPO starts from that adapter, samples eight
answers per training question, and runs 123 optimization steps. The training
judge assigns the equally weighted PASS fraction under the task rubric;
empty or truncated answers receive zero. This outcome reward differs from
the final reference-relative evaluation. Full recipes and provenance are
in Appendix~\ref{app:execution}.

\begin{table}[t]
\centering\small
\begin{tabular}{lrrr}
\toprule
Stage & No Skill & $\Delta_{\mathrm{Base}}$ & $\Delta_{\mathrm{prev}}$\\
\midrule
\inputtable{tables/internalization.tex}
\bottomrule
\end{tabular}
\caption{Reported 305-task No Skill aggregates for Qwen3-4B and its two
adapters. The latest per-task recovery archive remains to be reconciled
with this supplied aggregate version (Appendix~\ref{app:execution}).}
\label{tab:internalization}
\end{table}

The reported sequence is 41.56$\rightarrow$42.69$\rightarrow$43.04:
$+1.13$ points after SFT and another $+0.35$ after GRPO. This is a modest
descriptive improvement, not a demonstrated routing capability or a
significant effect. Training uses non-identical benchmark-derived variants,
so the comparison concerns related educational families rather than
unseen-family generalization. Updated adapter evaluations also recover
historical execution failures; item-level reconciliation is needed to
separate recovery and answer-quality contributions before finalizing the
learning claim.
